% 面向领导口头汇报：方法论精髓与实现效果（无封面，≤15 页）
% 编译：bash scripts/xelatex-clean.sh docs/reports/Agent预研形式方法-方法论精髓与效果-简报.tex
% 受众：数学 / 密码专家；不讲 PQC 入门与基础数论
\documentclass[UTF8,aspectratio=169,14pt]{ctexbeamer}
\usetheme{Madrid}
\usecolortheme{default}
\setbeamertemplate{navigation symbols}{}
\setbeamertemplate{footline}[frame number]
\setbeamerfont{frametitle}{size=\large,series=\bfseries}
\setbeamerfont{itemize/enumerate body}{size=\normalsize}
\usepackage{amsmath,amssymb,booktabs}
\usepackage{tikz}
\usetikzlibrary{arrows.meta,positioning}
\setlength{\parskip}{0.35em}

\begin{document}

% ---- 1 ----
\begin{frame}{问题：Agent 预研为何容易失控}
\begin{itemize}
  \item 开放生成空间过大：在文本里「找能跑通的拼法」，而非沿已验证能力生长。
  \item 「对拍过」被当成完成：缺依赖、禁源、接缝未认证，仍可假绿。
  \item 典型代价：同一 Encaps\,\texttt{correctness} 任务，公司/家里 Agent 各约一天，
        合计吃掉约 $2/3$ 月度配额（体感；非严格对账实验）。
\end{itemize}
\vspace{0.6em}
\begin{block}{要解决的不是「会不会写算子」}
而是：如何把开放生成收成\textbf{可检查、可拒绝}的受控搜索。
\end{block}
\end{frame}

% ---- 2 ----
\begin{frame}{精髓（一句话）}
\begin{center}
\large
只允许沿\\[0.4em]
\textbf{已验证能力及其合法组合}\\[0.4em]
生长实现；\\[0.6em]
\normalsize
缺项先补洞，禁止项不得当模板，\\
锁参遇阻则停——不许改参硬闯。
\end{center}
\vspace{0.8em}
\begin{itemize}
  \item 静态投影：依赖 DAG（谁依赖谁）。
  \item 合法拼装：哪些边可走、哪些在 $\mathrm{Forbidden}$。
  \item 作业过程：Agent 状态机上哪些转移被门禁允许。
\end{itemize}
\end{frame}

% ---- 3 ----
\begin{frame}{三层工具（给专家的压缩版）}
\begin{center}
\begin{tikzpicture}[
  box/.style={draw, rounded corners=2pt, align=center, font=\small,
              minimum width=3.6cm, minimum height=1.1cm},
  arr/.style={-{Stealth}, thick}
]
  \node[box, fill=blue!8] (g) {依赖图 $G=(V,E)$\\闭包 $\mathrm{Dep}^*(T)$};
  \node[box, fill=orange!10, right=8mm of g] (l) {合法拼装\\语义边 / 禁止边};
  \node[box, fill=green!10, right=8mm of l] (a) {作业过程\\允许的转移};
  \node[box, fill=gray!12, below=10mm of l] (ct) {闭包表 $\mathrm{CT}$：先交表，再写码};
  \draw[arr] (g) -- (ct);
  \draw[arr] (l) -- (ct);
  \draw[arr] (a) -- (ct);
\end{tikzpicture}
\end{center}
\vspace{0.4em}
\begin{itemize}
  \item 不在此展开完备性证明；工程上以 $\mathrm{CT}$ 为可审计接口。
  \item 证书阶梯：探针 $\to$ lemma $\to$ 交付；脏证书不可再当获证引用。
\end{itemize}
\end{frame}

% ---- 4 ----
\begin{frame}{闭包、缺项与「何时允许写码」}
对目标 $T$、已获证集 $\Gamma$：
\[
  D=\mathrm{Dep}^*(T),\qquad
  M=\mathrm{Missing}_\Gamma(T)=D\setminus\Gamma.
\]
\begin{block}{写码许可（操作判定）}
\begin{enumerate}
  \item $\mathrm{CT}$ 已交且与计划一致；
  \item $M=\emptyset$（含接缝项已获证或显式列入补洞计划且未宣称交付）；
  \item 计划与 $\mathrm{Forbidden}$ 相交为空；
  \item I/O 契约与验收级别 $\ell$ 已声明。
\end{enumerate}
\end{block}
缺任一项 $\Rightarrow$ \textbf{拒绝}「开始写交付代码 / 宣称完成」。
\end{frame}

% ---- 5 ----
\begin{frame}{Forbidden：负例与纪律}
\begin{columns}[T]
\begin{column}{0.48\textwidth}
\textbf{归入 Forbidden 的典型来源}
\begin{itemize}
  \item 已判决关闭的路线（\texttt{frozen/}）
  \item 「只求对、不计结构」的 correctness 标本
  \item 零垫凑维、limbsplit 等否决形状
\end{itemize}
\end{column}
\begin{column}{0.48\textwidth}
\textbf{允许的用法}
\begin{itemize}
  \item 读判决书：为何关闭、继任是谁
  \item 作正确性下界 / 过程对照
\end{itemize}
\vspace{0.4em}
\textbf{禁止}
\begin{itemize}
  \item 抄实现、抄 customspec、当模板 fork
\end{itemize}
\end{column}
\end{columns}
\vspace{0.8em}
\begin{alertblock}{一句话}
出门不带码——负例存在是为了挡住非法边，不是为了省事复用。
\end{alertblock}
\end{frame}

% ---- 6 ----
\begin{frame}{落到工程：两张「先交的表」}
\begin{enumerate}
  \item \textbf{闭包表 $\mathrm{CT}$}（接缝任务，如 Decaps）\\
        目标、依赖闭包、缺项、接缝、Forbidden 相交、验收级别。
  \item \textbf{参数卡 + 波次表}（参数组迁移，如 ML-KEM-768）\\
        锁 $(k,d_u,d_v)$、I/O 长度、分核/tiling；波次门禁：积木 $\to$ device $\to$ incubating。
\end{enumerate}
\vspace{0.5em}
\begin{itemize}
  \item 已锁形状 / \texttt{blockDim} / 分核：遇阻\textbf{停}并重开参数讨论。
  \item 每波固定：文书 $\to$ 实现 $\to$ CPU + SIM 双绿 $\to$ 再往下。
  \item golden 只验 I/O 等价；禁止「与参考源码同构」当验收。
\end{itemize}
\end{frame}

% ---- 7 ----
\begin{frame}{效果①：Decaps 试金石（接缝）}
\begin{itemize}
  \item \textbf{做法}：先交 $\mathrm{CT}_{\mathrm{decaps}}$，再按缺项序补洞、写码、验收（含拒绝路径）。
  \item \textbf{相对 correctness 探索（臂 A）}：强成功——搜索/返工形态收束；禁抄 Frozen/correctness。
  \item \textbf{相对人工优化交付主线（臂 C）}：正确性同阶；SIM tick 同量级；\textbf{不}宣称性能胜出。
\end{itemize}
\vspace{0.5em}
\begin{block}{判决}
先交表这道门\textbf{有用}；它压缩无效派生，\textbf{代替不了}人工优化主线，也堵不死验收脚本上的洞（须单独修假绿）。
\end{block}
\end{frame}

% ---- 8 ----
\begin{frame}{效果②：ML-KEM-768（参数组迁移）}
\begin{itemize}
  \item 压力点换到：真 $k=3$（奇数维）、禁零垫凑 $4/8$、全套 I/O 重锁。
  \item 路径：参数卡锁定 $\to$ W0--W1 积木 $\to$ W2--W3 device(+CT) $\to$ W4 incubating $\to$ glue。
  \item 关键不变量写死：polyvec6；Inner AIV\,$2{+}1$；INTT batch4；$\hat A[9]$ 分片 $5{+}4$。
\end{itemize}
\vspace{0.4em}
\begin{block}{完成口径（有条件）}
探针与 incubating（PKE$\times$3 + KEM$\times$3 + decaps-ct）CPU+SIM 绿；\\
AscendC-only KEM roundtrip 绿。\\
\textbf{未做}：stable-768、liboqs-768 交叉 KAT、NPU。
\end{block}
\end{frame}

% ---- 9 ----
\begin{frame}{效果数字（登记，非优化 SLA）}
\begin{columns}[T]
\begin{column}{0.52\textwidth}
\textbf{768 device SIM tick}（Cloud）
\begin{itemize}
  \item KeyGen $\approx 3.7\times 10^5$
  \item Encrypt $\approx 5.1\times 10^5$
  \item Decrypt $\approx 2.2\times 10^5$
  \item KEM Decaps $\approx 8.2\times 10^5$
\end{itemize}
相对同结构 $k{=}4$ 约低 $18\%$--$31\%$（维数预期内）。
\end{column}
\begin{column}{0.44\textwidth}
\textbf{正确性}
\begin{itemize}
  \item I/O 对拍 $\max=0$
  \item 拒绝路径：$K=J(z\Vert c)$且 $\neq$ accept
  \item KEM 闭环绿
\end{itemize}
\vspace{0.6em}
\small 详数见仓库\\
\texttt{qa/active\_sim\_regress\_summary.md}
\end{column}
\end{columns}
\end{frame}

% ---- 10 ----
\begin{frame}{成本体感（题外话，非倍数宣称）}
\begin{center}
\begin{tabular}{@{}lcc@{}}
  \toprule
  & 开放 correctness 探索 & 有门禁的 768 一夜 \\
  \midrule
  墙钟 & 两侧 Agent 各 $\sim 1$ 天 & $<4$ 小时 \\
  配额 & $\sim 2/3$ 月度 & $41\%\to 55\%$（$\Delta\sim 14$\,pp） \\
  终点形态 & 易得「能对」的膨胀结构 & incubating 全链可审计 \\
  \bottomrule
\end{tabular}
\end{center}
\vspace{0.7em}
\begin{itemize}
  \item 差的不是模型突然变强，而是\textbf{作业面被收小}。
  \item 数字为订阅面板与当场印象：序数级体感；\textbf{不}据此写「快 $N$ 倍」。
\end{itemize}
\end{frame}

% ---- 11 ----
\begin{frame}{证明了什么 / 没有证明什么}
\begin{columns}[T]
\begin{column}{0.48\textwidth}
\textbf{支持的判断}
\begin{itemize}
  \item 先交 $\mathrm{CT}$/参数卡可压缩无效派生
  \item Forbidden 纪律可执行
  \item 参数组迁移可按波次复用模式（非抄冻结码）
  \item Agent 可在门禁下走到可对拍实现
\end{itemize}
\end{column}
\begin{column}{0.48\textwidth}
\textbf{明确未声称}
\begin{itemize}
  \item 密码学机器验证证明
  \item 性能优于人工主线
  \item 已可无审合入 Rule/Skill
  \item 768 已交付 stable / 交叉 KAT
\end{itemize}
\end{column}
\end{columns}
\vspace{0.8em}
\begin{block}{地位}
调研草稿级方法论 + 两条压力测试床上的工程证据；不是最终理论封闭。
\end{block}
\end{frame}

% ---- 12 ----
\begin{frame}{可带走的三条}
\begin{enumerate}
  \item \textbf{把开放生成改成受控搜索}：只沿 $\Gamma$ 的合法闭包生长；$M\neq\emptyset$ 不许宣称交付。
  \item \textbf{负例与锁参同样重要}：Forbidden 出门不带码；已锁几何遇阻则停，不改参硬闯。
  \item \textbf{效果看过程与边界}：相对 correctness 强在轨迹；相对优化主线不靠 tick 吹牛；
        768 证明「换参数组」也可复用同一套门禁。
\end{enumerate}
\end{frame}

% ---- 13 ----
\begin{frame}{若需追问：材料指针}
\begin{itemize}
  \item 教材长文：\texttt{docs/research/从已验证能力到合法派生-…教材草案.pdf}\\
        （第 6--7 章 1024 复盘/前瞻；第 8 章 768 完整复盘 + 配额题外话）
  \item 768 参数卡：\texttt{docs/specs/fips203-mlkem768-parameter-card.md}
  \item 绿线与 tick：\texttt{qa/active\_sim\_regress\_summary.md}
  \item 当日纪要：\texttt{qa/2026-07/2026-07-27-768收尾复盘与文档刷新.md}
\end{itemize}
\vspace{0.8em}
\begin{center}
\small\textit{本简报无封面，共 13 页（含本页）。}
\end{center}
\end{frame}

\end{document}
